% glossary.tex - Ordliste og forkortelser for HDO VideoAPI

% ============================================================================
% FORKORTELSER - kun domenespesifikke og ikke-opplagte
% ============================================================================

% Helsevesen
\newacronym{hdo}{HDO}{Helsetjenestens driftsorganisasjon for nødnett HF}
\newacronym{amk}{AMK}{Akuttmedisinsk kommunikasjonssentral}
\newacronym{lvs}{LVS}{Legevaktsentral}
\newacronym{nhn}{NHN}{Norsk Helsenett}

% Tekniske
\newacronym{jwt}{JWT}{JSON Web Token}
\newacronym{oauth2}{OAuth 2.0}{Open Authorization 2.0}
\newacronym{webrtc}{WebRTC}{Web Real-Time Communication}
\newacronym{api}{API}{Application Programming Interface}
\newacronym{rest}{REST}{Representational State Transfer}
\newacronym{cicd}{CI/CD}{Continuous Integration/ Continuous Delivery}
\newacronym{vod}{VoD}{Video on Demand}
\newacronym{ki}{KI}{Kunstig intelligens}
\newacronym{http}{HTTP}{Hypertext Transfer Protocol}
\newacronym{uri}{URI}{Uniform Resource Identifier}
\newacronym{url}{URL}{Uniform Resource Locator}
\newacronym{dto}{DTO}{Data Transfer Object}
\newacronym{ice}{ICE}{Interactive Connectivity Establishment}
\newacronym{xml}{XML}{eXtensible Markup Language}
\newacronym{poc}{POC}{Proof-Of-Concept}
\newacronym{e2e}{E2E}{End-to-End}
\newacronym{pin}{PIN}{Personal Identification Number}
\newacronym{stun}{STUN}{Session Traversal Utilities for NAT}
\newacronym{turn}{TURN}{Traversal Using Relay around NAT}
\newacronym{cors}{CORS}{Cross-Origin Resource Sharing}
\newacronym{json}{JSON}{JavaScript Object Notation}
\newacronym{m2m}{M2M}{Machine-to-Machine}
\newacronym{di}{DI}{Dependency Injection}
\newacronym{hls}{HLS}{HTTP Live Streaming}
\newacronym{rtmp}{RTMP}{Real-Time Messaging Protocol}
\newacronym{vm}{VM}{Virtual Machine}
\newacronym{sdk}{SDK}{Software Development Kit}
\newacronym{ide}{IDE}{Integrated Development Environment}
\newacronym{tls}{TLS}{Transport Layer Security}
\newacronym{tcp}{TCP}{Transmission Control Protocol}

% ============================================================================
% FAGTERMER - kun prosjektspesifikke
% ============================================================================

\newglossaryentry{rest-api}{
    name={REST-API},
    description={API basert på Representational State Transfer (REST), der ressurser aksesseres via standard HTTP-metoder og unike URL-er}
}

\newglossaryentry{keycloak}{
    name={Keycloak},
    description={Åpen kildekode-løsning for identitets- og tilgangsstyring som 
    støtter OAuth 2.0 og OpenID Connect. Benyttes som 
    autorisasjonsserver for utstedelse av tilgangstokener i prosjektet}
}

\newglossaryentry{prosess-intern-state}{
    name={prosess-intern state},
    description={Data som kun lever i applikasjonsprosessens arbeidsminne og ikke persisteres til ekstern lagring, og derfor går tapt ved omstart av tjenesten}
}

\newglossaryentry{demo-frontend}{
    name={demo-frontend},
    description={En forenklet webapplikasjon utviklet for å demonstrere og 
    teste funksjonaliteten til en tjeneste. Er ikke ment for produksjonsbruk, 
    men fungerer som et visuelt grensesnitt mot et API}
}

\newglossaryentry{teknologisk-innlåsing}{
    name={teknologisk innlåsing},
    description={Situasjon der et system er så tett koblet til én leverandørs 
    teknologi, protokoler eller API-er at bytte av leverandør medfører 
    betydelige tekniske eller økonomiske kostnader}
}

\newglossaryentry{proof-of-concept}{
    name={proof-of-concept},
    description={En tidlig, forenklet implementasjon som har som mål å demonstrere at en teknisk løsning er gjennomførbar eller hensiktsmessig, uten å være produksjonsklar}
}

\newglossaryentry{provider-basert}{
    name={provider-basert},
    description={En arkitektur der funksjonalitet legges til som selvstendige, utskiftbare moduler (providers) uten å endre kjernesystemet}
}

\newglossaryentry{videomotor}{
    name={videomotor},
    description={Backend-system som håndterer videostrømming. I dette prosjektet: AntMedia Server og NHN Video}
}

\newglossaryentry{mellomvare}{
    name={mellomvare},
    description={Programvarelag som befinner seg mellom to systemer og håndterer kommunikasjon, oversettelse eller felles funksjonalitet. I denne rapporten: laget som eksponerer ett felles API mot klientene og videresender forespørsler til riktig videomotor}
}

\newglossaryentry{provider-laget}{
    name={provider-laget},
    description={Arkitekturlag som håndterer integrasjon mot eksterne videomotorer. Hvert provider-lag implementerer et felles grensesnitt, slik at kjernesystemet er uavhengig av underliggende videoteknologi}
}

\newglossaryentry{controller-laget}{
    name={controller-laget},
    description={Arkitekturlag som eksponerer API-endepunkter og håndterer innkommende HTTP-forespørsler. Delegerer forretningslogikk videre til underliggende lag}
}

\newglossaryentry{provider}{
    name={provider},
    description={Komponent som representerer rollen som integrasjonspunkt mot én spesifikk videomotor, implementert via \texttt{IVideoProvider}-grensesnittet}
}

\newglossaryentry{provider-adapter}{
    name={provider-adapter},
    description={Konkret implementasjon av en provider som oversetter mellom prosjektet sitt interne grensesnitt og en spesifikk videomotors API og dataformat}
}

\newglossaryentry{broadcast-service}{
    name={Broadcast Service},
    description={AntMedia-tjeneste for administrasjon av direktestrømmer og 
    videorom. Eksponerer endepunkter for oppretting, administrasjon og 
    avslutning av strømmer}
}

\newglossaryentry{vod-service}{
    name={\gls{vod} Service},
    description={AntMedia-tjeneste som håndterer forhåndsinnspilt videoinnhold. 
    Eksponerer endepunkter for opplasting, administrasjon og avspilling av 
    lagrede videoressurser}
}

\newglossaryentry{management-service}{
    name={Management Service},
    description={AntMedia-tjeneste for serverkonfigurasjon og systemadministrasjon. 
    Eksponerer endepunkter for administrasjon av serverinnstillinger og 
    driftsparametere}
}

\newglossaryentry{sha256}{
    name={HMAC-SHA256},
    description={Kryptografisk meldingsautentiseringskode som kombinerer en 
    hemmelig nøkkel med SHA-256-hashfunksjonen for å verifisere både 
    dataintegritet og avsenderautentisitet}
}

\newglossaryentry{password-grant}{
    name={password grant},
    description={OAuth 2.0-flyt der klienten sender brukernavn og passord 
    direkte til autorisasjonsserveren og mottar et tilgangstoken i retur}
}

\newglossaryentry{token}{
    name={token},
    description={Tidsbegrenset dataobjekt som utstedes av en autorisasjonsserver 
    og benyttes som bevis på at innehaveren har fått tilgang til en bestemt 
    ressurs eller tjeneste}
}

\newglossaryentry{sekvensiell}{
    name={sekvensiell utviklingsmodell},
    description={Utviklingsmodell der prosjektfasene gjennomføres i fast 
    rekkefølge, én etter én, uten tilbakevending til tidligere faser}
}

\newglossaryentry{fossefallsmodellen}{
    name={fossefallsmodellen},
    description={Sekvensiell utviklingsmodell der krav fastlegges ved oppstart 
    og prosjektet gjennomføres fase for fase — fra kravspesifikasjon via 
    design og implementasjon til testing og levering — uten iterasjon}
}

\newglossaryentry{iterativ}{
    name={iterativ arbeidsform},
    description={Utviklingstilnærming der arbeidet deles inn i kortere sykluser, 
    der hver syklus produserer kjørbar kode og gir grunnlag for justering av 
    krav og prioriteringer før neste syklus}
}

\newglossaryentry{scrum}{
    name={scrum},
    description={Agilt rammeverk for iterativ programvareutvikling organisert 
    rundt kortvarige sprinter, faste seremonier (planlegging, daglig stand-up, 
    retrospektiv) og løpende backlog-prioritering \cite{scrumguide}}
}

\newglossaryentry{git}{
    name={Git},
    description={Distribuert versjonskontrollsystem der hver utvikler har en 
    fullstendig kopi av prosjekthistorikken lokalt. Støtter parallell utvikling 
    via branching og brukes i dette prosjektet for kildekodehåndtering og 
    samarbeid via GitHub}
}

\newglossaryentry{pull-request}{
    name={pull-request},
    description={Forespørsel om å flette en ferdig utviklet kodegren inn i 
    en annen. Brukes som kontrollpunkt for kodegjennomgang før endringer 
    aksepteres i prosjektets felles kodebase}
}

\newglossaryentry{ci-pipeline}{
    name={CI-pipeline},
    description={Automatisert sekvens av steg som kjøres ved endringer i 
    kodebasen, typisk bygging, testing og distribusjon. I dette prosjektet 
    brukt via GitHub Actions for å verifisere at koden kompilerer og at 
    tester passerer ved hver pull request}
}

\newglossaryentry{regresjonsfeil}{
    name={regresjonsfeil},
    description={Feil som oppstår i tidligere fungerende funksjonalitet som 
    følge av nye kodeendringer. Avdekkes typisk av automatiserte tester som 
    kjøres ved hver endring i kodebasen}
}

\newglossaryentry{grenbasert}{
    name={grenbasert arbeidsflyt},
    description={Versjonskontrollpraksis der hver oppgave eller funksjonalitet 
    utvikles i en egen gren, isolert fra hovedgrenen}
}

\newglossaryentry{LL-HLS}{
    name={LL-HLS},
    description={Low-Latency HTTP Live Streaming — en utvidelse av HLS-protokollen 
    utviklet av Apple for å redusere forsinkelse i videostrømming til 2-5 sekunder}
}

\newglossaryentry{LL-DASH}{
    name={LL-DASH},
    description={Low-Latency Dynamic Adaptive Streaming over HTTP — en utvidelse 
    av DASH-protokollen med redusert forsinkelse, standardisert av MPEG}
}

\newglossaryentry{miljøvar}{
    name={miljøvariabel},
    description={Navngitt variabel definert i kjøretidsmiljøet utenfor 
    applikasjonskoden, brukt til å styre konfigurasjon som API-nøkler, 
    tilkoblingsstrenger og tjenesteparametere uten å hardkode verdier i kildekoden}
}

\newglossaryentry{constraints}{
    name={constraints},
    description={Arkitektoniske begrensninger som et system må overholde for å oppfylle en bestemt arkitekturstil}
}